\chapter*{Agradecimentos}
\addcontentsline{toc}{chapter}{Agradecimentos}

Aristóteles recorda que “a excelência é fruto do hábito”, destacando que o progresso é resultado de disciplina e constância. Marco Aurélio, por sua vez, ensina que “o que importa não é o que acontece, mas como reagimos ao que acontece”, enquanto Sêneca afirma que “dificuldades fortalecem a mente, assim como o trabalho molda o corpo”. Inspirado por essas perspectivas, reconheço que esta dissertação não representa apenas um produto acadêmico, mas um processo de formação intelectual e pessoal construído com perseverança — e, sobretudo, com apoio coletivo.

Agradeço profundamente aos meus colegas de mestrado, cuja dedicação, competência e espírito de cooperação elevaram, a cada etapa, o nível de exigência e reflexão. Os debates, desafios e contribuições mútuas criaram um ambiente acadêmico rico e estimulante. Como afirmou Kant, “o homem não se aperfeiçoa isoladamente”. Foi justamente pelo convívio com pares tão preparados que a régua se manteve sempre alta, impulsionando-me a aprimorar ideias, métodos e a mim mesmo.

Aos professores da FGV–EESP, registro meu sincero reconhecimento pelo rigor científico, clareza metodológica e compromisso contínuo com a excelência. Suas orientações e provocações intelectuais foram essenciais para a maturação deste trabalho. Em consonância com a visão nietzschiana de constante “superação de si mesmo”, seus ensinamentos ampliaram minha capacidade analítica e fortaleceram meu senso crítico.

Agradeço, em primeiro lugar, à minha família, cujo apoio incondicional, paciência e presença constante ofereceram o equilíbrio necessário nos momentos de maior intensidade. Aos meus pais e à minha esposa, pela compreensão, incentivo e serenidade que sustentaram cada etapa deste percurso. À minha filha, cuja luz e espontaneidade renovaram continuamente meu sentido de propósito. Como escreveu Nietzsche, “quem tem um porquê enfrenta quase qualquer como” — e minha família foi, em todos os momentos, esse “porquê”.

A todos que contribuíram direta ou indiretamente para esta jornada — colegas, professores, amigos e familiares — deixo registrado meu profundo agradecimento. Se esta dissertação alcança algum nível de excelência, é porque trilhei o caminho acompanhado por pessoas que me incentivaram a elevar meus hábitos, meu pensamento e minha atuação.